%==============================================================================
% SUPPORTING INFORMATION FOR JOURNAL OF CHEMICAL PHYSICS
%==============================================================================
\documentclass[aip,jcp,reprint]{revtex4-2}
\usepackage[utf8]{inputenc}
\usepackage[T1]{fontenc}
\usepackage{amsmath,amssymb,amsfonts,graphicx,booktabs,siunitx,physics,bm}
\usepackage{xcolor}
\usepackage{hyperref}

% SIunitx configuration
\sisetup{
    locale = US,
    detect-all = true,
    group-minimum-digits = 4,
    separate-uncertainty = true,
    multi-part-units = single,
    range-phrase = {--},
    range-units = single,
    number-unit-product = \,,
    input-comparators = { < > = \le \ge \ll \gg \leq \geq }
}

\begin{document}

\title{\textbf{Supporting Information for:\\Suppression of Pure Dephasing and Enhancement of Excitonic Coherence via Non-Markovian Spectral Bath Engineering}}

\author{Steve Cabrel Teguia Kouam}
\affiliation{Department of Physics, Faculty of Science, University of Douala, Cameroon}

\author{Theodore Goumai Vedekoi}
\affiliation{Department of Physics, Faculty of Science, University of Yaoundé I, Cameroon}

\author{Jean-Pierre Tchapet Njafa}
\affiliation{Department of Physics, Faculty of Science, University of Yaoundé I, Cameroon}

\author{Jean-Pierre Nguenang}
\affiliation{Department of Physics, Faculty of Science, University of Douala, Cameroon}

\author{Serge Guy Nana Engo}
\affiliation{Department of Physics, Faculty of Science, University of Yaoundé I, Cameroon}

\date{\today}
\maketitle

%==============================================================================
% SECTION S1: FMO HAMILTONIAN PARAMETERS
%==============================================================================
\section{FMO Hamiltonian and Spectral Density Parameterization}

The analytical details of the Process Tensor Hierarchy of Pure States (PT-HOPS), the Padé Matsubara decompositions, and the full 12-test numerical validation benchmarks isolating numerical positivity and Hierarchical Equations of Motion (HEOM) deviations are strictly documented inside the primary manuscript text.

This Supporting Information supplements those analytical proofs by explicitly reproducing the exact 7-site bacteriochlorophyll-\textit{a} matrices (obtained from Adolphs and Renger\cite{Adolphs2006}) utilized within our systemic calculations to assist independent computational replication vectors.

\subsection{Site Energies}

\begin{table}[h]
\centering
\caption{\textbf{FMO Site Energies (cm$^{-1}$)}}
\label{tab:SI_site_energies}
\begin{tabular}{ccccccc}
\toprule
$\varepsilon_1$ & $\varepsilon_2$ & $\varepsilon_3$ & $\varepsilon_4$ & $\varepsilon_5$ & $\varepsilon_6$ & $\varepsilon_7$ \\
\midrule
12400 & 12550 & 12250 & 12350 & 12450 & 12600 & 12500 \\
\bottomrule
\end{tabular}
\end{table}

\subsection{Electronic Couplings}
The inter-chromophore electronic coupling matrix $J_{nm}$ (in cm$^{-1}$):
\begin{equation}
\mathbf{J} = \begin{pmatrix}
0 & -87.7 & 5.5 & -5.9 & 6.7 & -13.7 & -9.9 \\
-87.7 & 0 & 30.8 & 8.2 & 0.7 & -11.9 & 4.3 \\
5.5 & 30.8 & 0 & -53.5 & -2.2 & -9.6 & 6.0 \\
-5.9 & 8.2 & -53.5 & 0 & 1.2 & -6.3 & 1.7 \\
6.7 & 0.7 & -2.2 & 1.2 & 0 & 9.6 & -6.0 \\
-13.7 & -11.9 & -9.6 & -6.3 & 9.6 & 0 & 89.7 \\
-9.9 & 4.3 & 6.0 & 1.7 & -6.0 & 89.7 & 0
\end{pmatrix}.
\label{eq:SI_coupling_matrix}
\end{equation}

\subsection{Spectral Density Parameters}
The underlying continuous Drude-Lorentz environment is mapped alongside four discrete vibronic bounds to precisely model non-Markovian oscillations. 
\begin{table}[h]
\centering
\caption{\textbf{Systematic Spectral Density Operating Parameters}}
\label{tab:SI_spectral_params}
\begin{tabular}{lc}
\toprule
\textbf{Parameter Constraint} & \textbf{Value Bound} \\
\midrule
Drude--Lorentz reorganization ($\lambda_{\rm DL}$) & \SI{35}{\per\centi\meter} \\
Drude--Lorentz thermal cutoff ($\gamma_{\rm DL}$) & \SI{50}{\per\centi\meter} \\
Underdamped Vibronic mode 1 ($\omega_1$) & \SI{150}{\per\centi\meter} \\
Underdamped Vibronic mode 2 ($\omega_2$) & \SI{200}{\per\centi\meter} \\
Underdamped Vibronic mode 3 ($\omega_3$) & \SI{575}{\per\centi\meter} \\
Underdamped Vibronic mode 4 ($\omega_4$) & \SI{1185}{\per\centi\meter} \\
Huang--Rhys coefficient $S_1$ & 0.05 \\
Huang--Rhys coefficient $S_2$ & 0.02 \\
Huang--Rhys coefficient $S_3$ & 0.01 \\
Huang--Rhys coefficient $S_4$ & 0.005 \\
\bottomrule
\end{tabular}
\end{table}

%==============================================================================
% REFERENCES
%==============================================================================
\begin{thebibliography}{99}
\bibitem{Adolphs2006}
Adolphs, J.; Renger, T. Biophys. J. \textbf{2006}, \textit{91}, 2778--2797.
\end{thebibliography}
\end{document}
